\begin{claim}{F.23}{Proved}{The based source-sign symmetry, complete source braid map, and two exact finite marginal images}
Retain the original groups, source map, basepoints, and two
one-point forgettings of F.19--F.22. In particular put
\[
 \mathcal K_\pm=p_\pm(K^+)\leq R_\pm\leq V_\pm,
       \qquad V_\pm\cong\widehat F_7,
                                                               \tag{F.23.1}
\]
with the actual geometric positive coordinates and Artin negative
coordinates of F.22. The source involution
$h(s)=\zeta^2/s$, together with the explicit path $\lambda$ below,
gives an actual based isomorphism of inclusions
\[
 \begin{array}{ccc}
 \mathcal K_-&\lhook\joinrel\longrightarrow&R_-\\
 {\scriptstyle\mathcal B_+}\downarrow\phantom{\scriptstyle\sim}
       &&\phantom{\scriptstyle\sim}\downarrow{\scriptstyle\mathcal B_+}\\
 \mathcal K_+&\lhook\joinrel\longrightarrow&R_+ .
 \end{array}                                                  \tag{F.23.2}
\]
The same map extends to $V_-\xrightarrow{\sim}V_+$.
Its formula, inverse, and all filling squares are specified below.

There is a complete exact braid description of the actual homomorphism
\[
 \Phi^\dagger:L^\dagger=\widehat{\pi_1(Y^\dagger,s_0)}
                       =\widehat F_{27}\longrightarrow G.
                                                               \tag{F.23.3}
\]
It uses the twenty-seven specified source lassos, a rational
projection perturbation, the original endpoint whiskers, and the
original F.20 path $\eta$. The exact arithmetic certificates and
twenty-seven full braid words are retained in the accompanying
files. The coordinate comparison is the clockwise-exchange and
inverse triangular comparison of F.20.7--9 and F.22.11--13.

For each sign, for each of the two groups $H=S_3,A_4$, and for
every continuous homomorphism from the indicated full point-pushing
group, the following equality holds:
\[
 \rho(\mathcal K_\pm)=\rho(R_\pm)
       \qquad\bigl(\rho\in\operatorname{Hom}_{\rm cont}(V_\pm,H),
                              \ H\in\{S_3,A_4\}\bigr).
                                                               \tag{F.23.4}
\]
This is the exact finite scope of the exhaustive calculations.
The original maps $\mathcal K_\pm\hookrightarrow R_\pm$ and
$K^+\hookrightarrow D^+$ remain the specified inclusions;
their surjectivity remains unresolved. The joint image and its
kernel remain exactly those of F.22.8a.
\end{claim}

\begin{proof}
\emph{The literal source and the sign-changing maps.}
Write
\[
 \begin{gathered}
 a(s)=\frac{\zeta-s}{\zeta+s},\qquad
 v_u(s)=\left(\frac{u-s}{u+s}\right)^2,\qquad
 u\in\{1,\zeta,\zeta^2,\zeta^3,\zeta^4\},\\
 A_u=u(2\zeta-u),\qquad
 B_u=\frac{\zeta u^2}{2u-\zeta}
       \quad(u\ne\zeta),\\
 \Delta_+=\{s:s^2=A_u,\ u\ne\zeta\},\qquad
 \Delta_-=\{s:s^2=B_u,\ u\ne\zeta\},\\
 Y^\dagger=Y\setminus(\Delta_+\cup\Delta_-).
 \end{gathered}                                                \tag{F.23.5}
\]
The complete finite configuration is
\[
 \phi^\dagger(s)=
 (0,1,v_1(s),v_\zeta(s),v_{\zeta^2}(s),
                     v_{\zeta^3}(s),v_{\zeta^4}(s),a(s),-a(s)),
                                                               \tag{F.23.6}
\]
with the original additional label at infinity. Its normalization
and its induced map to $G$ are those of F.19 and F.22.14.
All punctures and their simple collision multiplicities are
the actual ones established in F.17.

For $h(s)=\zeta^2/s$ and $u'=\zeta^2/u$, direct substitution gives
\[
 \begin{gathered}
 h^2(s)=s,\qquad a(h(s))=-a(s),\qquad
 v_u(h(s))=v_{u'}(s),\\
 \frac{\zeta^4}{A_u}=B_{u'},\qquad
 \frac{\zeta^4}{B_u}=A_{u'}.
 \end{gathered}                                                \tag{F.23.7}
\]
Thus $h$ is an actual homeomorphism of $Y^\dagger$ and $Y$,
and it interchanges the two signed sources
$Y_\pm=Y\setminus\Delta_\pm$.
Let $\Pi$ permute the old labels by $u\mapsto u'$ and exchange
the two new labels. It fixes $0,1,\infty,v_\zeta$.
The complete source equality is
\[
                  \phi^\dagger h=\Pi\phi^\dagger .            \tag{F.23.8}
\]
This equality commutes with all the original fillings and with
the two one-point forgettings, with their signs exchanged.

The based path is explicit. Put $s_1=h(s_0)$ and use the polygon
\[
                 \lambda=(s_0, s_0+6i, s_1+6i, s_1).
                                                               \tag{F.23.9}
\]
It lies in $Y^\dagger$. The exact sign-symmetry producer checks
that each of its two vertical lines has real coordinate different
from every old puncture. For either real coordinate $x$ and every
collision square $V=A_u,B_u$, it checks $x\ne0$ and
\[
                  4x^4-4\operatorname{Re}(V)x^2
                            -(\operatorname{Im}V)^2\ne0.     \tag{F.23.10}
\]
Indeed a root $x+iy$ of $z^2=V$ would have
$y=\operatorname{Im}(V)/(2x)$ and would make this expression zero.
The middle segment has imaginary part greater than two, whereas
every old and collision puncture has absolute value less than two.
These are exact inequalities in the retained real number field.

With left-to-right path products, the actual based source
automorphism and its inverse are
\[
 \begin{aligned}
 \mathcal A_{\rm src}(g)&=\lambda h(g)\lambda^{-1},\\
 \mathcal A_{\rm src}^{-1}(g)&=h(\lambda)^{-1}h(g)h(\lambda).
 \end{aligned}                                                \tag{F.23.11}
\]
Both composites cancel to $g$. The square
$\mathcal A_{\rm src}^2$ is conjugation with the literal
loop $\lambda h(\lambda)$.
These maps extend to inverse continuous maps of the full
source completions and commute with the original filling $q$.

Put $w=\phi^\dagger(\lambda)$, a path from the original target
basepoint $c_0$ to $\Pi(c_0)$. The actual target automorphism is
\[
 \mathcal A_G(z)=w\Pi(z)w^{-1},\qquad
 \mathcal A_G^{-1}(z)=\Pi^{-1}(w^{-1}zw).                       \tag{F.23.12}
\]
The literal equality (F.23.8) gives
$\mathcal A_G\Phi^\dagger=\Phi^\dagger\mathcal A_{\rm src}$.
It follows on the actual source filling kernel that
$\mathcal A_G(K^+)=K^+$; the commuting target fillings give
$\mathcal A_G(D^+)=D^+$ and preserve the corresponding old-point
kernel $J$.

Write $F_+$ and $F_-$ for the full $G$ forgettings in F.22.
Let $\Pi_+$ be the old-label permutation together with renaming
the retained negative new label as positive. Then
$F_+\Pi=\Pi_+F_-$. Set $w_+=F_+(w)$. The two inverse maps are
\[
 \begin{aligned}
 \mathcal B_+(r)&=w_+\Pi_+(r)w_+^{-1},\\
 \mathcal B_+^{-1}(r)&=\Pi_+^{-1}(w_+^{-1}rw_+),\qquad
 F_+\mathcal A_G=\mathcal B_+F_- .
 \end{aligned}                                                \tag{F.23.13}
\]
The old-point filling square identifies its restrictions
$V_-\to V_+$, and the common-three-point filling square identifies
$R_-\to R_+$. These are isomorphisms with the displayed inverse.
Finally the equality on $K^+$ and the last formula give
$\mathcal B_+(\mathcal K_-)=\mathcal K_+$, proving (F.23.2).

\emph{A full geometric source basis.}
Use $s_*=2+2i$ and axis-parallel squares of half-side
$h_0=10^{-4}$. Their twenty-seven centers are $0$, both signs of
the five old centers in the first table, and both signs of the
eight centers in the second table. Every pair $(r,t)$ below
denotes $(r+it)/10^6$.
\[
 \begin{array}{c|rr}
 \text{old label}&r&t\\\hline
 1&1000000&0\\
 \zeta&309017&951057\\
 \zeta^2&-809017&587785\\
 \zeta^3&-809017&-587785\\
 \zeta^4&309017&-951057
 \end{array}
 \qquad
 \begin{array}{cc|rr}
 u&\text{square}&r&t\\\hline
 1&A_u&882644&1077509\\
 1&B_u&454951&555393\\
 \zeta^2&A_u&80730&-1390528\\
 \zeta^2&B_u&41611&-716735\\
 \zeta^3&A_u&1260728&-1131556\\
 \zeta^3&B_u&439300&-394290\\
 \zeta^4&A_u&1685063&174411\\
 \zeta^4&B_u&587159&60773
 \end{array}                                                  \tag{F.23.14}
\]
For each square take the straight whisker from $s_*$ to its
northeast corner, traverse its boundary counterclockwise, and
return along the same whisker. Call these loops $\gamma_0,\ldots,
\gamma_{26}$ in the producer's recorded order: the eleven old
punctures first, then the sixteen signed collisions.

Here are the exact checks which make these actual free generators.
For every old center the radius-$h_0/2$ disk contains its named
old root. For a new center $c$ and its named square $V$, the checks give
\[
 |c|>\tfrac12,\qquad |c^2-V|<10^{-5},\qquad
 (h_0/2)^2+10^{-5}<h_0/2.                                    \tag{F.23.15}
\]
On the circle of radius $h_0/2$ about $c$, the linear term
$2c(z-c)$ has absolute value greater than $h_0/2$, while
$(z-c)^2+c^2-V$ has smaller absolute value. Rouch\'e's theorem
therefore gives exactly one root of $z^2-V$ in that disk.
The twenty-seven squares are pairwise disjoint by the exact
center-distance inequalities. Each whisker stays farther than
$\sqrt2h_0$ from every other square's center. Its basepoint is
northeast of its own square, and its first intersection is that
square's northeast corner. All whisker directions are distinct.
They consequently form an embedded common-base star. Cutting
the punctured plane along this star gives the usual free basis
of its fundamental group. Filling the sixteen collision punctures
retains precisely the eleven old free generators.

\emph{The exact projection perturbation and crossing certificate.}
The field is the same as in F.20:
\[
 \mathbb Q(\theta),\qquad
 \theta^4-5\theta^2+5=0,\qquad 19/10<\theta<2.
\]
Real signs are decided by rational isolation of this indicated
root and exact field arithmetic. Put $P_0=1+371i/1000$ and use
the literal offsets
\[
 \delta_j=\frac{10^{-14}j^2}{P_0},
                  \qquad 0\leq j\leq8,                       \tag{F.23.16}
\]
in the finite-label order of (F.23.6).
For each segment and each pair of strands write their original
difference as $N/D$. The producer verifies the polynomial
inequality, over the entire real parameter interval,
\[
 |P_0N|^2-\bigl(10^{-14}(i^2-j^2)\bigr)^2|D|^2>0.
                                                               \tag{F.23.17}
\]
Every denominator is nonzero on the specified segment.
This strict inequality proves that the entire interpolation
$z_j\mapsto z_j+t\delta_j$, $0\leq t\leq1$, is collision-free.
The polynomial is certified by rational Bernstein subdivision
with strictly positive coefficients on every terminal interval.
All its coefficients and terminal intervals are retained.

For the perturbed paths, real projected differences are again
exact field polynomials after multiplying by positive squared
denominators. Bernstein--Descartes subdivision isolates all their
zeros, and the certificate checks nonzero height and derivative
on each root interval. These intervals are refined until all
crossing times are disjoint. At each crossing the two labels
are adjacent in the current projection order; the height sign
then gives the clockwise Artin letter calibrated in F.20.
The final order is checked against the exact endpoint order.
The record retains the full polynomials, root intervals, height
and derivative signs, subdivision counts, and resulting words.
These operations prove the full path word by successive actual
projection-chamber comparisons.

The primary run has $137$ forward segments, $4,932$ pair checks,
and $1,412$ crossings. Its exact sign algorithm made $846,769$
field-sign calls. Return whiskers are the literal reverse paths.
The additional chart path below has five segments and ten
crossings. Thus the complete comparison uses $142$ forward
segments, $5,112$ pair checks, and $1,422$ crossings. The initial
unperturbed projection had simultaneous events from $0,a,-a$;
its rejected run is retained separately. The proved interpolation
(F.23.17) is the precise comparison used for the completed run.

\emph{The standard chart and the original endpoint whiskers.}
At $s_*$ let $b=\phi^\dagger(s_*)$. Let $r$ have the same old
coordinates as $b$ and the two new coordinates
$R_p=8+3i$, $R_n=12+6i$. The actual path $T:r\to b$ first
moves $R_p$ straight to $a(s_*)$ with $R_n$ fixed, and then
moves $R_n$ straight to $-a(s_*)$ with the positive point fixed.
Both paths and their entire offset interpolations are certified.

The standard configuration is
\[
 r_{\rm std}=(0,1,3,2,4,5,6,8+3i,12+6i).                      \tag{F.23.18}
\]
The path $G:r_{\rm std}\to r$ moves the five old nonconstant
labels one at a time along the specified straight segments,
keeping both new points and $0,1$ fixed. The certified projected
order of $r_{\rm std}$ is the F.22 order
$0,1,v_\zeta,v_1,v_{\zeta^2},v_{\zeta^3},v_{\zeta^4},p,n$.
For an explicit straightening, write these standard projected
points as $x_j+iy_j$ and take the continuous piecewise linear
function $f$ with $f(x_j)=y_j$, affine between adjacent $x_j$
and constant outside them. The oriented plane map
\[
 x+iy\longmapsto x+i(y-f(x))
                                                               \tag{F.23.19}
\]
has inverse $x+iy\mapsto x+i(y+f(x))$. Its interpolation using
$tf$ changes no projected order. This gives the actual chart
and inverse used for the clockwise generators and the lower
whiskers; the same chart is restricted for every filling.

At every endpoint $z$ retain the actual path
$H_z(t)=(z_j+t\delta_j)_j$. The basepath producer checks
the original and shifted orders of all $284$ segment endpoints,
in $10,224$ exact pair comparisons. Every $H_z$ stays in its
original projection-order chamber. Writing a subscript $1$
for an offset path, the certified two-parameter homotopies give
the following identities in the actual path groupoid:
\[
 \begin{gathered}
 \gamma=H_b\gamma_1H_b^{-1},\qquad
 TH_b=H_rT_1,\qquad GH_r=H_{r_{\rm std}}G_1,\\
 GT\gamma T^{-1}G^{-1}
 =H_{r_{\rm std}}
       (G_1T_1\gamma_1T_1^{-1}G_1^{-1})H_{r_{\rm std}}^{-1}.
 \end{gathered}                                                \tag{F.23.20}
\]
Thus the computed word is carried to the original standard
configuration by its displayed endpoint whisker. The chamber
comparison preserves precisely the named clockwise generators.
The first two coordinates on an offset loop are the constant
values $0,1+t\delta_1$. Dividing all coordinates by
$1+t\delta_1$, which is nonzero by the same pairwise certificate,
is the literal affine normalization of F.22.14. Consequently
(F.23.20) also holds in every original normalized target.

Retain the original F.20.4a path $\eta:s_0\to s_*$.
Its existing exact certificate places it in $Y^\dagger$:
the radial first segment misses the old and inner collision
roots by its radius, and misses all outer roots by the four
nonzero imaginary parts $\operatorname{Im}(A_u\overline{s_0^2})$.
The remaining segments lie outside radius two.
Let $\mathcal N$ be the exact affine normalization of F.22.14.
The full original comparison path in the normalized target is
\[
             \Xi=\mathcal N(GT\phi^\dagger(\eta)^{-1}).         \tag{F.23.21}
\]
All three paths $G,T,\phi^\dagger(\eta)$ keep their first
two coordinates literally $0,1$. Hence the normalization
$(z_j-z_0)/(z_1-z_0)$ leaves every remaining coordinate
unchanged at every point of this entire path. Thus $\Xi$
has exactly the original standard and $s_0$ endpoints.
Let $b_j$ denote the recorded standard nine-strand braid
of $GT\phi^\dagger(\gamma_j)T^{-1}G^{-1}$, interpreted through
(F.23.20). Then the complete source formula is
\[
 \Phi^\dagger(\eta\gamma_j\eta^{-1})
           =\Xi^{-1}\mathcal N_*(b_j)\Xi,
                          \qquad0\leq j\leq26.                 \tag{F.23.22}
\]
This follows by cancelling the identical based paths; it keeps
the original source basepoint and every original target label.

Extend the finite-point path $\Xi$ by successive supported local
disk isotopies. The image of the standard disk containing the
first eight points and excluding the ninth specifies the
permitted F.22 disk at the original base configuration. The
inverse isotopy and (F.23.19) give its inverse chart. Applying
the isotopy to a loop gives a square whose two side paths are
$\Xi$; that square identifies the induced fiber-coordinate map
with the groupoid conjugation in (F.23.22). The original F.22
section still has the constant comparison paths at $(p,n)$.
Forgetting labels from the same square proves all its based
filling comparisons.

The twenty-seven actual source generators freely generate
$\pi_1(Y^\dagger,s_0)$. Their recorded images therefore give
its entire homomorphism. The universal full completion maps
give the unique continuous map on $L^\dagger$ with these images.
It equals the original completed $\Phi^\dagger$, since the
two maps agree on its dense discrete source. This proves
the complete-source assertion (F.23.3).

\emph{The two retained source sections and the exact marginal maps.}
The eleven old generators give a particular section
$i_{\rm geo}:L\to L^\dagger$ of the original filling $q$.
Retain the F.19 section $i$ as well. On the actual target
coordinates put
\[
 c(\ell)=\Phi^\dagger(i_{\rm geo}(\ell)i(\ell)^{-1})\in K^+,
       \qquad j_{\rm geo}(\ell)=c(\ell)j(\ell).                \tag{F.23.23}
\]
These formulas follow because both sections have the same
$q$-image $\ell$. The exact comparison of the two pullback
semidirect coordinates is
\[
 (d,\ell)\longmapsto(d c(\ell)^{-1},\ell),
       \qquad(e,\ell)\longmapsto(e c(\ell),\ell).              \tag{F.23.24}
\]
Both maps express the same actual element, and their
composites are the identity. The original F.19 meridians
keep their original whiskers. The twenty-seven new lassos
and the old generators define the same actual source
filling kernel; restricting the same $\Phi^\dagger$ to
that kernel gives exactly $K^+$ in either coordinate system.

Delete strand $9$ from each $b_j$ to obtain its positive
eight-strand braid, delete strand $8$ to obtain its negative
one, and delete both to obtain the seven-old-strand braid.
The algorithm tracks the literal labels through every crossing.
The resulting paths and (F.23.22) are the full $G$ forgetting
maps $F_+,F_-$ before their restrictions to $J$.
The inverse triangular fiber maps of F.20 and F.22 commute
with these fillings: a deleted lower-whiskered clockwise
generator becomes trivial, and every retained one becomes
the corresponding generator in the retained order.
Equivalently, filling the same variables in
\[
 \tau_i(X)=X_n^{-1}\cdots X_{i+1}^{-1}
                         X_i^{-1}X_{i+1}\cdots X_n
                                                               \tag{F.23.25}
\]
gives exactly the triangular word in that retained order.
The additional calibration producer checks the full pure-braid
filling identity on every standard pure generator at the ranks
and deletion sets used here, as well as the Artin relations and
both triangular inverse maps. Thus its finite list of free-word
identities proves those homomorphism identities on the entire
pure braid groups.

Each of the sixteen collision words has trivial old braid and
trivial opposite-sign marginal braid, both by free reduction
of the recorded words and by the actual forgetting maps.
Its remaining linking vector is supported on its selected
old/new pair with value $-1$. This is the counterclockwise
simple holomorphic collision in the retained clockwise Artin
convention. The full words, including this sign, are retained.

\emph{Exact finite evaluation of the marginal maps.}
The following recurrence evaluates each recorded eight-strand
word in any finite group, without expanding intermediate free
words. Begin with $x_1,\ldots,x_7$, set $x_8=1$, and set all
conjugator values $c_j=1$. For a positive letter $\sigma_j$ use
\[
 (x_j,x_{j+1},c_j,c_{j+1})
  \longmapsto
 (x_jx_{j+1}x_j^{-1},\,x_j,\,x_jc_{j+1},\,c_j),                 \tag{F.23.26}
\]
and for its inverse use
\[
 (x_j,x_{j+1},c_j,c_{j+1})
  \longmapsto
 (x_{j+1},\,x_{j+1}^{-1}x_jx_{j+1},\,c_{j+1},
                                      \,x_{j+1}^{-1}c_j).     \tag{F.23.27}
\]
Induction on the displayed Artin substitutions proves these
identities for the images and their labelled conjugators at
each prefix. At a pure final word put $u=c_8^{-1}$; the final
seven $x_i$ evaluate the filled automorphism $\alpha$.
The actual conjugation on the fiber, proved by the cocycle
and inverse triangular comparison, evaluates as
\[
                         X_i\longmapsto u x_i u^{-1}.          \tag{F.23.28}
\]
For the negative marginal these are precisely the Artin
coordinates of $V_-$. For the positive marginal, a seven-tuple
of values on its actual geometric generators first undergoes
$\tau_7$; use those as the initial $x_i$. A collision seed then
has value $u$. For a source action apply $\tau_7$ also to the
final tuple in (F.23.28). This is exactly the evaluation of
the two inverse coordinate maps, since $\tau_7^2=1$.

There are eight collision seeds for either sign, and eleven
old source automorphisms. They are the computed actual
images and actions for $i_{\rm geo}$. The image of the signed
source filling kernel is the closed subgroup generated by
all their $L$-conjugates, as follows from the free source basis
and its full completed filling. Its image is exactly
$\mathcal K_\pm$ by F.22.21 and (F.23.23)--(F.23.24).

\emph{The finite candidate reduction with actual peripheral maps.}
F.21 and the actual surjections $p_\pm:D^+\to R_\pm$ give
\[
        \overline{\langle\!\langle\mathcal K_\pm
                    \rangle\!\rangle}_{R_\pm}=R_\pm.          \tag{F.23.29}
\]
Consequently in every finite image its image normally
generates the image of $R_\pm$, and has that entire image
when the latter is nilpotent, by F.21's finite argument.

The four deleted-puncture coordinates also retain their
actual meridians. For a collision over old label $u$, the
simple local parameter $a-v_u$, or $-a-v_u$, takes its
counterclockwise source circle once around that collision.
Transport its whisker by the labelled plane isotopy along
the specified source path. The inverse endpoint isotopy
carries the local circle to a based peripheral loop in the
original seven-punctured plane. Comparing its transported
whisker with the specified lower whisker gives the actual
conjugating loop between these two peripheral loops.
In the positive geometric coordinates this identifies the
seed with a conjugate of $u_j^{-1}$. In the negative
coordinates (F.23.25) identifies $U_j$ with a conjugate of
the inverse clockwise lasso; the seed is consequently a
conjugate of $U_j$. These comparisons preserve orientation
and the named puncture. In particular, in every finite
homomorphism the seed and its named deleted coordinate
have exactly the same order.

For $H=S_3$, the proper subgroups are the trivial subgroup,
three subgroups $C_2$, and the normal subgroup $C_3$.
For $H=A_4$, they are the trivial subgroup, three subgroups
$C_2$, four subgroups $C_3$, and the normal subgroup $V_4$.
The $S_3$ list follows directly from its three transpositions
and two three-cycles; a subgroup with two distinct
transpositions, or with both a transposition and a
three-cycle, is the full group. The $A_4$ list, its actual
element sets and normal closures, is also obtained by the
exhaustive finite subgroup closure algorithm in the
coefficient certificate. All these
proper subgroups are nilpotent. The proper subgroups
normally generating $S_3$ are its three $C_2$ subgroups;
those normally generating $A_4$ are its four $C_3$ subgroups.

Thus a failure of (F.23.4) requires both
$\rho(V_\pm)=H$ and $\rho(R_\pm)=H$. Its marginal image
then lies in a proper subgroup $M=C_2$ for $S_3$, or
$M=C_3$ for $A_4$. Conjugation in $H$ carries these respective
subgroups to one fixed $M$, so all homomorphisms with such
a failure occur in the finite enumeration with that fixed
subgroup. The peripheral comparison just proved makes all
four deleted-coordinate values have order one or two in
the first case, and order one or three in the second.
At least one is nonidentity, since they normally generate
the image of $R_\pm$.

Every seven-tuple in $H^7$ is the values of a unique
continuous homomorphism from the actual $\widehat F_7$.
The complete possible-failure candidate counts are therefore
\[
 \begin{aligned}
 S_3:&\quad 6^3(4^4-1)-3\,2^3(2^4-1)=54,720,\\
 A_4:&\quad 12^3(9^4-1)-4\,3^3(3^4-1)=11,327,040.
 \end{aligned}                                                \tag{F.23.30}
\]
The subtracted families are exactly the tuples contained
in a proper subgroup while having a nonidentity deleted
coordinate. They are disjoint because that coordinate
determines the particular $C_2$ or $C_3$ subgroup.

\emph{The exhaustive source-orbit test and its precise conclusion.}
For each candidate tuple evaluate the eight actual collision
words by (F.23.26)--(F.23.28), and retain the tuple precisely
when all eight values lie in the fixed $M$. Call the finite
set of retained tuples $\mathcal T_0$. Each of the eleven
actual old source automorphisms acts as a permutation of
$\operatorname{Hom}(V_\pm,H)$, by precomposition.
It preserves generation of the image. Its actual pure
labelled point-pushing action carries each deleted
peripheral generator to a conjugate of itself, so it
also preserves the candidate order restrictions and the
condition that some deleted value is nonidentity.
Record all eleven images of every element of $\mathcal T_0$.
Delete an element when one of these images is outside
$\mathcal T_0$, and repeat on the retained set.

This computes exactly the largest subset of
$\mathcal T_0$ invariant under the eleven permutations.
Indeed a subset mapped into itself by a permutation of a
finite set is also invariant under its inverse. Hence
forward invariance under the eleven maps is invariance
under their generated group. Their profinite source
completion has the same finite orbits: the continuous
action on this finite representation set has dense
discrete source image and therefore equal finite image.
The resulting invariant set consists exactly of the
candidates for which every source-conjugate collision
seed has image in $M$. Those source-conjugates generate
the actual marginal subgroup, so this is the exact
remaining failure set.

Both signs were enumerated separately. The complete
executed results are
\[
 \begin{array}{c|r|r|r|r}
 H&\text{candidates}&|\mathcal T_0|&
       \text{first deletion}&\text{survivors}\\\hline
 S_3&54,720&624&624&0\\
 A_4&11,327,040&15,552&15,552&0
 \end{array}                                                  \tag{F.23.31}
\]
The table holds for each sign. The outputs retain every
element of $\mathcal T_0$, its seven coordinates, all
eleven image keys and edges, and its deletion layer.
Multiplication is literal composition of the enumerated
permutations. The $A_4$ run streams the same exhaustive
candidate set in finite batches; it does not sample it.
The empty invariant sets prove (F.23.4) for every specified
continuous homomorphism, including the cases already
settled by the finite nilpotent image argument.

\emph{Executed certificate and retained scope.}
The complete source geometry is produced by
\texttt{full\_cross\_source\_geometry\_stage.py};
\texttt{full\_cross\_source\_basepaths\_stage.py}
retains the chart and endpoint comparisons;
\texttt{full\_cross\_marginal\_braids\_stage.py}
records all full and deleted words.
The two finite producers are
\texttt{marginal\_s3\_source\_stage.py} and
\texttt{marginal\_a4\_source\_stage.py}.
The sign-symmetry producer passes $64$ exact checks,
and \texttt{source\_artin\_calibration\_stage.py}
passes $2,473$ exact polynomial and free-word checks.
They are retained under \texttt{code/part\_a/} with their
full outputs, stdout, and hashes. Independent isolated
execution reproduced all $35$ complete-source, symmetry,
and finite-scan output files byte for byte; the calibration
output was independently reproduced as well.

All original sections, fillings, signs, and based
comparison paths enter the displayed maps. The finite
claim is the equality of the two images under every map
$V_\pm\to S_3$ or $V_\pm\to A_4$. The remaining full
marginal inclusion and the interaction subgroup
$K^+\cap W$ of F.22 retain their original source-membership
questions.
\end{proof}
